\section{תורת הפרעות מנוונת}
\date{01.01.2026}
\\
נבחן מילטוניאן מהצורה
\[
H = H_0 + \lambda V,
\]
כאשר $\lambda \ll 1$ הוא פרמטר קטן חסר ממדים, ו$H_0$ ניתן לפתרון מלא.

נניח כי ל$H_0$ קיימת רמת אנרגיה מנוונת
\[
H_0 \ket{n,s^{(0)}} = E_n^{(0)} \ket{n,s^{(0)}},
\qquad s = 1,\dots,d,
\]
כאשר $d$ הוא ממד תת־המרחב המנוון. כל יתר המצבים העצמיים של $H_0$ יסומנו על־ידי $\ket{k^{(0)}}$ עם אנרגיות $E_k^{(0)} \neq E_n^{(0)}$.

נפתח את פונקציית הגל והאנרגיה בטור חזקות ב$\lambda$:
\begin{align}
\ket{n,s} &= \ket{n,s^{(0)}} + \lambda \ket{n,s^{(1)}} + \lambda^2 \ket{n,s^{(2)}} + \cdots, \\
E_{n,s} &= E_n^{(0)} + \lambda E_{n,s}^{(1)} + \lambda^2 E_{n,s}^{(2)} + \cdots.
\end{align}

הסימון העליון $(k)$ מציין \emph{סדר בפיתוח בלמדה}, ואינו קשור כלל לערכים עצמיים של $H_0$.

\subsection{משוואות הסדרים}

הצבת הפיתוחים במשוואת שרדינגר
\[
H\ket{n,s} = E_{n,s}\ket{n,s}
\]
ואיסוף חזקות של $\lambda$ נותנים:

\paragraph{סדר אפס}
\[
H_0 \ket{n,s^{(0)}} = E_n^{(0)} \ket{n,s^{(0)}}.
\]

\paragraph{סדר ראשון}
\[
H_0 \ket{n,s^{(1)}} + V \ket{n,s^{(0)}}
=
E_n^{(0)} \ket{n,s^{(1)}} + E_{n,s}^{(1)} \ket{n,s^{(0)}}.
\]

\paragraph{סדר שני}
\[
H_0 \ket{n,s^{(2)}} + V \ket{n,s^{(1)}}
=
E_n^{(0)} \ket{n,s^{(2)}}
+ E_{n,s}^{(1)} \ket{n,s^{(1)}}
+ E_{n,s}^{(2)} \ket{n,s^{(0)}}.
\]

וכן הלאה.



\subsection{הצעד הקריטי: לכסון ההפרעה בתת־המרחב המנוון}

מגדירים את מטריצת ההפרעה בתת־המרחב המנוון:
\[
V_{sr} = \mel{n,s^{(0)}}{V}{n,r^{(0)}}.
\]

\begin{resultbox}
\\
הצעד הראשון וההכרחי בתורת הפרעות מנוונת הוא ללכסן את $V_{sr}$.
\end{resultbox}

נבחר בסיס חדש $\{\ket{n,s^{(0)}}\}$ כך ש
\[
V_{sr} = \delta_{sr} V_{ss}.
\]

בבסיס זה נקבל מייד מהמשוואה מסדר ראשון, על־ידי כפל משמאל ב$\bra{n,r^{(0)}}$:
\[
E_{n,s}^{(1)} = V_{ss}.
\]



\subsection{חישוב התיקון מסדר ראשון לפונקציית הגל}

נכתוב את $\ket{n,s^{(1)}}$ בפיתוח מלא בבסיס שלם של $H_0$:
\[
\ket{n,s^{(1)}} =
\sum_{r=1}^d \ket{n,r^{(0)}} \braket{n,r^{(0)}}{n,s^{(1)}}
+
\sum_{k\notin n} \ket{k^{(0)}} \braket{k^{(0)}}{n,s^{(1)}}.
\]

אין כל סיבה עקרונית שהאיברים $\braket{n,r^{(0)}}{n,s^{(1)}}$ יתאפסו.

כדי לחשב את המקדמים מחוץ לתת־המרחב המנוון, נכפיל את משוואת הסדר הראשון משמאל ב$\bra{k^{(0)}}$:
\[
\mel{k^{(0)}}{H_0}{n,s^{(1)}} + \mel{k^{(0)}}{V}{n,s^{(0)}}
=
E_n^{(0)} \braket{k^{(0)}}{n,s^{(1)}}.
\]

מאחר ש
\[
\mel{k^{(0)}}{H_0}{n,s^{(1)}} = E_k^{(0)} \braket{k^{(0)}}{n,s^{(1)}},
\]
נקבל
\[
(E_k^{(0)} - E_n^{(0)}) \braket{k^{(0)}}{n,s^{(1)}}
=
-\mel{k^{(0)}}{V}{n,s^{(0)}}.
\]

ולכן
\[
\boxed{
\braket{k^{(0)}}{n,s^{(1)}} =
\frac{\mel{k^{(0)}}{V}{n,s^{(0)}}}{E_n^{(0)} - E_k^{(0)}}
}.
\]

המכנה לעולם אינו מתאפס, משום ש$k$ שייך מחוץ לתת־המרחב המנוון.



\subsection{קביעת הרכיבים בתוך תת־המרחב – סדר שני}

כדי לקבוע את המקדמים $\braket{n,r^{(0)}}{n,s^{(1)}}$ יש לעבור לסדר הבא.

נכפיל את משוואת הסדר השני משמאל ב$\bra{n,r^{(0)}}$ עם $r \neq s$:
\[
\mel{n,r^{(0)}}{V}{n,s^{(1)}}
=
E_{n,s}^{(1)} \braket{n,r^{(0)}}{n,s^{(1)}}.
\]

נציב את הפיתוח של $\ket{n,s^{(1)}}$ ונקבל:
\[
V_{rr} \braket{n,r^{(0)}}{n,s^{(1)}}
+
\sum_{k\notin n}
\frac{
\mel{n,r^{(0)}}{V}{k^{(0)}}
\mel{k^{(0)}}{V}{n,s^{(0)}}
}{
E_n^{(0)} - E_k^{(0)}
}
=
V_{ss} \braket{n,r^{(0)}}{n,s^{(1)}}.
\]

מאחר ש$V$ אלכסוני בתת־המרחב:
\[
(V_{ss} - V_{rr}) \braket{n,r^{(0)}}{n,s^{(1)}}
=
\sum_{k\notin n}
\frac{
\mel{n,r^{(0)}}{V}{k^{(0)}}
\mel{k^{(0)}}{V}{n,s^{(0)}}
}{
E_n^{(0)} - E_k^{(0)}
}
.
\]

אם $V_{ss} \neq V_{rr}$ נקבל
\[
\boxed{
\braket{n,r^{(0)}}{n,s^{(1)}}
=
\frac{
\displaystyle
\sum_{k\notin n}
\frac{
\mel{n,r^{(0)}}{V}{k^{(0)}}
\mel{k^{(0)}}{V}{n,s^{(0)}}
}{
E_n^{(0)} - E_k^{(0)}
}
}{
V_{ss} - V_{rr}
}
}.
\]



\subsection{ניוון שלא הוסר – לכסון היררכי}

אם מתקיים $V_{ss} = V_{rr}$, צד שמאל מתאפס אך צד ימין אינו חייב להתאפס – מתקבלת סתירה.

במקרה זה מגדירים מטריצה אפקטיבית מסדר שני:
\[
M_{rs}
=
\sum_{k\notin n}
\frac{
\mel{n,r^{(0)}}{V}{k^{(0)}}
\mel{k^{(0)}}{V}{n,s^{(0)}}
}{
E_n^{(0)} - E_k^{(0)}
},
\]
ומלכסנים אותה בתוך תת־המרחב שנותר מנוון.

\begin{resultbox}
\\
הסרת ניוון מתבצעת היררכית:  
אם ניוון אינו מוסר בסדר ראשון – עוברים לסדר שני, ואם צריך גם לשלישי.
\end{resultbox}



\section{אטום המימן – מבוא לפיצול הדק}

ההמילטוניאן הלא־יחסותי:
\[
H_0 = \frac{p^2}{2\mu} - \frac{e^2}{4\pi\varepsilon_0 r},
\]
עם רמות אנרגיה
\[
E_n^{(0)} = -\frac{1}{2}\mu c^2 \frac{\alpha^2}{n^2},
\qquad
\alpha = \frac{e^2}{4\pi\varepsilon_0 \hbar c}.
\]

הניוון הוא $2n^2$ (כולל ספין).



\section{ההמילטוניאן של הפיצול הדק}

התיקונים היחסותיים מסדר $\alpha^2$:
\[
H_{\mathrm{FS}}
=
-\frac{p^4}{8m^3c^2}
+
\frac{1}{2m^2c^2}\frac{1}{r}\frac{dV}{dr}\,\vec L\cdot\vec S
+
\frac{\hbar^2}{8m^2c^2}\nabla^2 V.
\]

שלושת האיברים הם:
\begin{itemize}
\item תיקון יחסותי לאנרגיה הקינטית,
\item צימוד ספין–תנע זוויתי,
\item איבר דרווין (מריחה קוונטית של הפוטנציאל).
\end{itemize}



\section{הערכת סדרי גודל}

מאחר ש
\[
\frac{p^2}{2m} \sim \alpha^2 mc^2,
\]
נקבל
\[
\frac{p^4}{8m^3c^2} \sim \alpha^4 mc^2.
\]

בדומה,
\[
\frac{1}{2m^2c^2}\frac{1}{r}\frac{dV}{dr} L\cdot S
\sim
\alpha^4 mc^2,
\qquad
\frac{\hbar^2}{8m^2c^2}\nabla^2 V
\sim
\alpha^4 mc^2.
\]

\begin{resultbox}
\\
הפיצול הדק הוא תיקון מסדר $\alpha^2$ לאנרגיות אטום המימן.
\end{resultbox}

מאחר ש$\alpha \simeq 1/137$, תורת ההפרעות מתכנסת במהירות רבה, והאלקטרומגנטיות היא תורה הפרעתית באופן עמוק.
\section*{תורת ההפרעות למצבים מנוונים והמבנה הדק}

\subsection*{1. השלמות לתורת ההפרעות המנוונת}
בשיעור הקודם עסקנו במתודולוגיה לטיפול במערכת עם ניוון ברמות האנרגיה. נשלים כעת את חישוב התיקון לפונקציית הגל מסדר ראשון ונתאר את האלגוריתם לטיפול במצבים בהם הניוון אינו מוסר לחלוטין בסדר ראשון.

\paragraph{1.1 הגדרת הבעיה}
נתונה מערכת עם המילטוניאן $H = H_0 + \lambda V$. נניח כי ל$H_0$ יש רמת אנרגיה $E_n^{(0)}$ המנוונת $g$ פעמים. נסמן את תתהמרחב המנוון ב$D$.
המצבים העצמיים בתתהמרחב יסומנו ב $|n^0, s\rangle$ כאשר $s=1,\dots,g$.
עבור מצבים מחוץ לתתהמרחב המנוון ($k \notin D$), האנרגיה היא $E_k^{(0)} \neq E_n^{(0)}$.

הפיתוח לטור חזקות של $\lambda$:
$$ E_{ns} = E_n^{(0)} + \lambda E_{ns}^{(1)} + \lambda^2 E_{ns}^{(2)} + \dots $$
$$ |n, s\rangle = |n^0, s\rangle + \lambda |n^1, s\rangle + \dots $$

\paragraph{1.2 תיקון מסדר ראשון לאנרגיה (חזרה)}
השלב הראשון והקריטי בתורת הפרעות מנוונת הוא לכסון מטריצת ההפרעה בתתהמרחב המנוון.

\begin{resultbox}
\\
\textbf{תיקונים}

התיקונים מסדר ראשון לאנרגיה, $E_{ns}^{(1)}$, הם הערכים העצמיים של המטריצה $V$ המוגבלת לתתהמרחב $D$:
\[ V_{rs} = \langle n^0, r | V | n^0, s \rangle \]
המצבים העצמיים של מטריצה זו מהווים את ה"בסיס הנכון" ($Good Basis$) לסדר אפס, $|n^0, s\rangle$, המקיים:
\[ \langle n^0, r | V | n^0, s \rangle = E_{ns}^{(1)} \delta_{rs} \]
\end{resultbox}

\paragraph{1.3 תיקון מסדר ראשון לפונקציית הגל}
אנו מחפשים את הביטוי ל$|n^1, s\rangle$. נפתח אותו בבסיס המצבים של $H_0$:
\begin{equation}
|n^1, s\rangle = \sum_{k \notin D} c_k |k^0\rangle + \sum_{r \in D} c_r |n^0, r\rangle
\end{equation}
את המקדמים $c_k$ עבור המצבים שאינם מנוונים ($k \notin D$) ניתן למצוא ממשוואת סדר ראשון (בדומה לתורה לא מנוונת):
\[ c_k = \frac{\langle k^0 | V | n^0, s \rangle}{E_n^{(0)} - E_k^{(0)}} \]
הבעיה היא למצוא את המקדמים $c_r$ בתוך תתהמרחב המנוון, שכן המכנה מתאפס. משוואת סדר ראשון אינה מספקת מידע עליהם. לשם כך, עלינו לפנות למשוואת סדר שני.

\paragraph{1.4 שימוש במשוואת סדר שני}
משוואת ההפרעה לסדר $k=2$ היא:
\[ H_0 |n^2, s\rangle + V |n^1, s\rangle = E_n^{(0)} |n^2, s\rangle + E_{ns}^{(1)} |n^1, s\rangle + E_{ns}^{(2)} |n^0, s\rangle \]
נבצע הטלה (Project) של המשוואה על מצב אחר בתתהמרחב המנוון, $\langle n^0, r|$ (כאשר $r \neq s$):
\[ \langle n^0, r | H_0 | n^2, s\rangle + \langle n^0, r | V | n^1, s\rangle = E_n^{(0)} \langle n^0, r | n^2, s\rangle + E_{ns}^{(1)} \langle n^0, r | n^1, s\rangle \]
מכיוון ש$H_0$ הרמיטי ופועל על מצב עצמי שמאלי, האיבר הראשון באגף שמאל מתבטל עם האיבר הראשון באגף ימין ($E_n^{(0)}$). נותרנו עם:
\[ \langle n^0, r | V | n^1, s\rangle = E_{ns}^{(1)} \langle n^0, r | n^1, s\rangle \quad (\text{עבור } r \neq s) \]
נציב את הפיתוח של $|n^1, s\rangle$ (משוואה 1) בתוך האיבר השמאלי:
\[ \sum_{k \notin D} \frac{\langle n^0, r | V | k^0 \rangle \langle k^0 | V | n^0, s \rangle}{E_n^{(0)} - E_k^{(0)}} + \sum_{r' \in D} c_{r'} \underbrace{\langle n^0, r | V | n^0, r' \rangle}_{V_{rr'} = E_{nr}^{(1)}\delta_{rr'}} = E_{ns}^{(1)} c_r \]
בגלל שבחרנו את הבסיס המלוכסן, הסכום השני קורס לאיבר יחיד $V_{rr} c_r$. נקבל:
\[ \sum_{k \notin D} \frac{V_{rk} V_{ks}}{E_n^{(0)} - E_k^{(0)}} + V_{rr} c_r = V_{ss} c_r \]
כאשר השתמשנו בסימון $V_{ss} = E_{ns}^{(1)}$. נעביר אגפים ונבודד את $c_r$ (המקדם המבוקש המייצג את הערבוב בתוך תתהמרחב):

\begin{resultbox}
\\
המקדם המגדיר את התרומה של המצב המנוון $|n^0, r\rangle$ לתיקון מסדר ראשון של המצב $|n^0, s\rangle$ הוא:
\[ \langle n^0, r | n^1, s\rangle = \frac{1}{V_{ss} - V_{rr}} \sum_{k \notin D} \frac{V_{rk} V_{ks}}{E_n^{(0)} - E_k^{(0)}} \]
\end{resultbox}

\paragraph{1.5 הסרת ניוון שיורי (סדר שני)}
מה קורה אם גם לאחר לכסון $V$, עדיין נותר ניוון? כלומר $V_{ss} = V_{rr}$?
במקרה כזה, המכנה בביטוי הנ"ל מתאפס והמשוואה אינה עקבית אלא אם המונה מתאפס גם הוא.
זהו מצב הדורש "לכסון מסדר שני". אנו מגדירים מטריצה אפקטיבית חדשה $M$ הפועלת בתתהמרחב שנותר מנוון:

\begin{definitionbox} 
\\
מטריצת הפרעה מסדר שני
עבור קבוצת המצבים בהם $E_{ns}^{(1)} = E_{nr}^{(1)}$, נגדיר:
\[ M_{rs} \equiv \sum_{k \notin D} \frac{V_{rk} V_{ks}}{E_n^{(0)} - E_k^{(0)}} \]
יש ללכסן את המטריצה $M$. הערכים העצמיים שלה יתנו את תיקוני האנרגיה מסדר שני, והווקטורים העצמיים יגדירו את בסיס האפס הנכון להמשך הפיתוח.
\end{definitionbox}



\subsection*{2. המבנה הדק באטום המימן (\LR{Fine Structure})}

לאחר ההצלחה של משוואת שרדינגר בשחזור ספקטרום אטום המימן (סדרת בלמר וכו'), מדידות מדויקות יותר הראו פיצולים נוספים ברמות האנרגיה. פיצולים אלו נובעים מאפקטים יחסותיים וספין.

\paragraph{2.1 ההמילטוניאן}
ההמילטוניאן השלם כולל את האיבר הקינטי, הפוטנציאל הקולוני, ותיקוני המבנה הדק:
\[ H = H_0 + H_{FS} \]
כאשר $H_0 = \frac{p^2}{2m} - \frac{e^2}{4\pi\epsilon_0 r}$ ותיקון המבנה הדק מורכב משלושה איברים:
\[ H_{FS} = H_{rel} + H_{SO} + H_{Darwin} \]

פירוט האיברים:
1.  תיקון יחסותי לאנרגיה הקינטית ($H_{rel}$):
    נובע מתוך פיתוח טיילור של האנרגיה היחסותית $E = \sqrt{p^2c^2 + m^2c^4}$. האיבר הראשון הוא המסה, השני הוא הקינטי הקלאסי, והשלישי הוא התיקון:
    \[ H_{rel} = -\frac{p^4}{8m^3c^2} \]
2.  צימוד ספיןמסילה ($H_{SO}$):
    נובע מהאינטראקציה בין המומנט המגנטי של האלקטרון (ספין) לשדה המגנטי שהוא "רואה" במערכת המנוחה שלו (בגלל תנועת הגרעין סביבו):
    \[ H_{SO} = \frac{e^2}{8\pi\epsilon_0} \frac{1}{m^2c^2 r^3} \vec{L} \cdot \vec{S} \]
3.  איבר דרווין ($H_{Darwin}$):
    אפקט קוונטייחסותי (קשור ל\LR{Zitterbewegung}) הנובע מכך שהאלקטרון אינו ממוקם נקודתית אלא "מרוח" על סקאלה של אורך גל קומפטון. רלוונטי רק למצבי $s$ ($l=0$):
    \[ H_{Darwin} = \frac{\hbar^2}{8m^2c^2} 4\pi \left(\frac{e^2}{4\pi\epsilon_0}\right) \delta^{(3)}(\vec{r}) \]

\paragraph{2.2 סדרי גודל וקבוע המבנה הדק}
הפרמטר הקטן בפיתוח הוא קבוע המבנה הדק:
\[ \alpha = \frac{e^2}{4\pi\epsilon_0 \hbar c} \approx \frac{1}{137} \]
באמצעות ניתוח ממדים, ניתן להראות כי כל שלושת האיברים הם מאותו סדר גודל:
\[ \Delta E_{FS} \sim \alpha^2 E_n^{(0)} \]
זהו תיקון של סדר גודל $10^{-4}$ ביחס לאנרגיה המקורית, מה שמצדיק שימוש בתורת ההפרעות.



\subsection*{3. חישוב התיקון לרמת היסוד ($n=1$)}

נסתכל על רמת היסוד של אטום המימן: $n=1, l=0$.
רמה זו מנוונת פעמיים בגלל הספין ($m_s = \pm 1/2$).
המצבים הם: $|1,0,0, \uparrow\rangle, |1,0,0, \downarrow\rangle$.

\paragraph{3.1 הסרת הניוון (או היעדרו)}
עלינו לבדוק את השפעת ההפרעה $H_{FS}$ על תתהמרחב המנוון.
1.  צימוד ספיןמסילה: מכיוון ש$l=0$, אופרטור התנע הזוויתי $\vec{L}$ נותן 0 על המצבים. לכן:
    \[ \langle H_{SO} \rangle_{n=1} = 0 \]
    האיבר הזה לא תורם ולא מסיר ניוון.
2.  תיקון יחסותי ודרווין: איברים אלו אינם תלויים בספין (הם סקלרים במרחב הספין). לכן המטריצה שלהם בתתהמרחב של הספין היא אלכסונית (פרופורציונלית למטריצת היחידה).

מסקנה: הניוון של הספין אינו מוסר על ידי המבנה הדק ברמת היסוד. שני מצבי הספין מקבלים את אותה הסטת אנרגיה.

\paragraph{3.2 חישוב הסטת האנרגיה}
נחשב את ערך התצפית (תיקון מסדר ראשון):
\[ E_{1}^{(1)} = \langle 1s | H_{rel} + H_{Darwin} | 1s \rangle \]
לאחר ביצוע האינטגרלים המרחביים (שימוש בפונקציות הגל של אטום המימן ובקשר $p^2 = 2m(EV)$), מתקבלת התוצאה הסופית עבור רמת היסוד:

\begin{resultbox}
\\
\textbf{אנרגיית רמת היסוד}

אנרגיית רמת היסוד כולל המבנה הדק
\[ E_1 \approx E_1^{(0)} \left( 1 + \frac{\alpha^2}{4} \right) \]
או בצורה מפורשת יותר:
\[ E_1 = -\frac{1}{2} m c^2 \alpha^2 \left( 1 + \frac{\alpha^2}{4} \right) \]
\end{resultbox}
הערה: הביטוי המלא לתיקון המבנה הדק עבור רמה כללית $n, j$ הוא:
\[ E_{nj}^{(1)} = \frac{E_n^{(0)} \alpha^2}{n^2} \left( \frac{n}{j+1/2} - \frac{3}{4} \right) \]